\section{Problem and Scope}

An agent self-improvement event is not merely another inference, tool call, or execution trace. Ordinary agent actions usually affect one task instance: a model answers a question, invokes a tool, retrieves context, or updates an intermediate working state. A self-improvement event is different because it changes an artifact that may shape later decisions, such as a prompt, policy, tool wrapper, memory snapshot, retrieval configuration, planner, evaluation script, or release candidate. This makes the accountability problem temporal. The relevant question is not only what happened in one run, but why a reusable change was proposed, which prior version it derived from, what evidence supported it, whether the candidate was validated, why a gate accepted or rejected it, and how the system can roll back if the change is later found unsafe or invalid.

Prior work and infrastructure address important parts of this space. Self-Refine, Reflexion, and Voyager show that language agents can improve behavior through feedback, reflection, or accumulated skills \cite{self_refine_2023,reflexion_2023,voyager_2023}. AgentDevel frames self-evolving agents as a release-engineering process involving traces, candidate versions, and regression gates \cite{agentdevel_2026}. Observability systems such as OpenTelemetry GenAI and LangSmith provide trace and feedback structures that can record execution-level evidence \cite{otel_genai_semconv,langsmith_observability_concepts}. Provenance and packaging frameworks such as W3C PROV, FDO, and Workflow Run RO-Crate provide reusable foundations for activities, entities, agents, digital-object records, and machine-actionable research objects \cite{w3c_prov_dm,fdo_architecture_spec,workflow_run_rocrate}.

What remains missing is a minimal, implementation-independent evidence profile for the self-improvement event itself. Such a profile should not replace self-improvement algorithms, tracing systems, provenance models, or packaging frameworks. Instead, it should define the self-improvement-specific constraints that connect these layers: which evidence is required before a candidate can be frozen, how references are bound to artifacts, how a validation result relates to a gate decision, and how rollback evidence remains attached to the promoted or rejected change. In this sense, ASIEP acts as a thin evidence-object layer between raw traces and reusable evidence packages.

The scope of this paper is deliberately narrow. We do not propose a new agent optimization algorithm, a benchmark for agent capability, or a production governance system. We also do not claim full OpenTelemetry integration, real LangSmith API integration, FDO registry registration, PID issuance, full RO-Crate certification, or external standard certification. We focus on a local minimal implementation of ASIEP, including a schema, lifecycle state machine, invariants, validator, repair planner, resolver, trace importer, local FDO/RO-Crate-like packager, evaluator, fixture corpus, and adversarial evidence cases.
